\chapter{Anexo C: modelos de volatilidad}

Este anexo sintetiza \path{doc/docs/models/}: vision general, features, splits, entrenamiento, reentrenamiento y familias.

\section{Vision general}

La parte de modelos toma \path{VOLATILITY_DB}, selecciona columnas, aplica split temporal train/validation/test, controla contratos compartidos, genera features, ejecuta k-folds temporales, selecciona hiperparametros, reentrena y evalua en test. La variable objetivo es \path{ImpliedVolatility}.

\section{Caracteristicas}

\path{features.md} distingue variables originales y derivadas. Las originales son tipo de opcion, strike, subyacente, tiempo hasta vencimiento y tipo. Las derivadas incluyen \path{TTEYears}, \path{sqrtTTEYears}, \path{logMoneyness}, \path{logMoneynessSq}, \path{logMoneynessXSqrtTTE}, \path{logForwardMoneyness} y \path{rate}. El dashboard conserva tambien variables raw para lectura financiera.

\section{Splits y k-folds}

\path{splits-kfolds.md} documenta split temporal, exclusividad de contratos, k-folds temporales, lookahead bias y data snooping. La prioridad es que test no contamine trainval y que validation no comparta contratos con train. Los k-folds se construyen dentro del bloque de entrenamiento.

\section{Entrenamiento}

\path{training.md} describe busqueda de hiperparametros, metricas base, metrica de seleccion, early stopping, escalado y artefactos. Las metricas base son MAE, RMSE y $R^2$. La seleccion compuesta penaliza error medio, variabilidad entre folds y brecha train-validation. Las redes y XGBoost usan validacion interna para early stopping.

\section{Reentrenamiento}

\path{retraining-progressive.md} documenta reentrenamiento no progresivo y progresivo. La fase final ajusta con trainval y evalua en test. El progressive training segmenta por moneyness para incorporar primero zonas mas cercanas a ATM y despues alas.

\section{Familias}

\path{families.md} compara regresion lineal, Random Forest, XGBoost, redes neuronales secuenciales y red tensor-train. \path{linear-regression.md} actua como baseline interpretable. \path{random-forest.md} capta no linealidades mediante ensembles. \path{xgboost.md} usa boosting regularizado. \path{neural-networks.md} describe redes densas con escalado y early stopping. \path{tensor-train.md} introduce una capa inspirada en tensor-train para interacciones compactas.
